\documentclass[11pt,a4paper]{article}

% Preambulo compartido por todas las partes del informe.
% Cada parte hace % Preambulo compartido por todas las partes del informe.
% Cada parte hace % Preambulo compartido por todas las partes del informe.
% Cada parte hace \input{preambulo} y despues carga sus propias tablas.

\usepackage[utf8]{inputenc}
\usepackage[T1]{fontenc}
\usepackage[spanish,es-noquoting]{babel}
\usepackage[a4paper,margin=20mm]{geometry}
\usepackage{amsmath,amssymb}
\usepackage{booktabs}
\usepackage{array}
\usepackage{xcolor}
\usepackage{enumitem}
\usepackage{titlesec}

% ------------------------------------------------ colores
% 'gris' lo usan las tablas generadas para las entradas repetidas por simetria
\definecolor{gris}{RGB}{170,170,170}

% ------------------------------------------------ espaciado
\titleformat{\section}{\normalfont\large\bfseries}{\thesection}{0.6em}{}
\titlespacing*{\section}{0pt}{2.4ex plus 1ex}{1.2ex}
\titleformat{\subsection}{\normalfont\bfseries}{\thesubsection}{0.6em}{}
\titlespacing*{\subsection}{0pt}{1.8ex plus .5ex}{0.8ex}

\setlength{\parindent}{0pt}
\setlength{\parskip}{0.6ex}
\renewcommand{\arraystretch}{1.08}

% ------------------------------------------------ atajos
\newcommand{\hextt}[1]{\texttt{#1}}
 y despues carga sus propias tablas.

\usepackage[utf8]{inputenc}
\usepackage[T1]{fontenc}
\usepackage[spanish,es-noquoting]{babel}
\usepackage[a4paper,margin=20mm]{geometry}
\usepackage{amsmath,amssymb}
\usepackage{booktabs}
\usepackage{array}
\usepackage{xcolor}
\usepackage{enumitem}
\usepackage{titlesec}

% ------------------------------------------------ colores
% 'gris' lo usan las tablas generadas para las entradas repetidas por simetria
\definecolor{gris}{RGB}{170,170,170}

% ------------------------------------------------ espaciado
\titleformat{\section}{\normalfont\large\bfseries}{\thesection}{0.6em}{}
\titlespacing*{\section}{0pt}{2.4ex plus 1ex}{1.2ex}
\titleformat{\subsection}{\normalfont\bfseries}{\thesubsection}{0.6em}{}
\titlespacing*{\subsection}{0pt}{1.8ex plus .5ex}{0.8ex}

\setlength{\parindent}{0pt}
\setlength{\parskip}{0.6ex}
\renewcommand{\arraystretch}{1.08}

% ------------------------------------------------ atajos
\newcommand{\hextt}[1]{\texttt{#1}}
 y despues carga sus propias tablas.

\usepackage[utf8]{inputenc}
\usepackage[T1]{fontenc}
\usepackage[spanish,es-noquoting]{babel}
\usepackage[a4paper,margin=20mm]{geometry}
\usepackage{amsmath,amssymb}
\usepackage{booktabs}
\usepackage{array}
\usepackage{xcolor}
\usepackage{enumitem}
\usepackage{titlesec}

% ------------------------------------------------ colores
% 'gris' lo usan las tablas generadas para las entradas repetidas por simetria
\definecolor{gris}{RGB}{170,170,170}

% ------------------------------------------------ espaciado
\titleformat{\section}{\normalfont\large\bfseries}{\thesection}{0.6em}{}
\titlespacing*{\section}{0pt}{2.4ex plus 1ex}{1.2ex}
\titleformat{\subsection}{\normalfont\bfseries}{\thesubsection}{0.6em}{}
\titlespacing*{\subsection}{0pt}{1.8ex plus .5ex}{0.8ex}

\setlength{\parindent}{0pt}
\setlength{\parskip}{0.6ex}
\renewcommand{\arraystretch}{1.08}

% ------------------------------------------------ atajos
\newcommand{\hextt}[1]{\texttt{#1}}

% generado por gen_parte_b.py -- no editar a mano
\newcommand{\dmin}{8}
\newcommand{\nCosetsCorregibles}{2325}
\newcommand{\nCosetsAmbiguos}{1771}


\begin{document}

\begin{center}
  {\large\bfseries Código de Golay extendido $(24,12)$ --- Parte B}\\[0.3ex]
  {\small Trabajo Práctico 2 · Curso de FEC 2026 · Fundación Fulgor}\\[0.3ex]
  {\small Villar, Pedro}
\end{center}

El modelo de referencia está en el paquete \texttt{golay/} y usa la
librería \texttt{fec\_algebra} del TP1. Las tablas de este informe las
genera \texttt{gen\_parte\_b.py} corriendo el modelo.

\section{Ejercicio 4 --- $GF(2^m)$ con $m=1$}

Se instancia el campo del TP1 con $m=1$ y $P(x)=x+1$:
\texttt{GF(m=1, primitive\_poly=0b1)}. Queda un campo de dos elementos
con estas tablas:

\begin{center}\small
\begin{tabular}[t]{@{}c|cc@{}}
$+$ & 0 & 1 \\
\midrule
0 & 0 & 1 \\
1 & 1 & 0 \\
\end{tabular}
\qquad
\begin{tabular}[t]{@{}c|cc@{}}
$\cdot$ & 0 & 1 \\
\midrule
0 & 0 & 0 \\
1 & 0 & 1 \\
\end{tabular}
\qquad
\begin{tabular}[t]{@{}c|c@{}}
$a$ & $a^{-1}$ \\
\midrule
0 & --- \\
1 & 1 \\
\end{tabular}
\end{center}

La suma es XOR y el producto es AND, o sea que es $GF(2)$. No hubo que
modificar la librería. El producto con $m=1$ hace una sola iteración y
da $a\wedge b$, y el inverso se calcula como $a^{2^m-2}=a^0=1$, que es
correcto porque el único elemento invertible es el 1. El 0 tira
\texttt{GFZeroDivisionError}.

Lo que faltaba era el álgebra de matrices, que está en
\texttt{golay/linalg.py}: \texttt{vec\_mat} (producto vector-matriz),
\texttt{mat\_mat} (producto matriz-matriz) y \texttt{weight} (peso de
Hamming), más algunas auxiliares (\texttt{transpose}, \texttt{identity},
\texttt{hstack}, \texttt{vec\_add}). Sirven para cualquier $GF(2^m)$, no
solo para $GF(2)$.

\section{Ejercicio 5 --- Construcción y verificación}

$B$ se carga con las doce filas del enunciado y el resto se arma con
\texttt{linalg}: $G$ = \texttt{hstack(I, B)}, $H$ = \texttt{hstack(B, I)}
y $H^T$ = \texttt{transpose(H)}. Con los productos matriciales se
verifica que $B^T=B$, $B^2=I$ y $G\,H^T=0$, igual que en la Parte A.

Codificando los 4096 mensajes con $v=m\cdot G$ y contando pesos:

\begin{center}\small
% generado por gen_parte_b.py -- no editar a mano
\begin{tabular}{@{}l*{5}{c}r@{}}
\toprule
peso $w$ & 0 & 8 & 12 & 16 & 24 & total \\
\midrule
modelo (fuerza bruta) & 1 & 759 & 2576 & 759 & 1 & 4096 \\
enunciado & 1 & 759 & 2576 & 759 & 1 & 4096 \\
\bottomrule
\end{tabular}

\end{center}

Coincide con la tabla del Ejercicio 2, y el peso mínimo no nulo da
$d_{\min}=\dmin$.

\section{Ejercicio 6 --- Codificador y decodificador}

El codificador es $v=m\cdot G$ con \texttt{vec\_mat}. Para
$m=\hextt{A5C}$ da $v=\hextt{A5C9A5}$, igual que a mano.

El decodificador tiene una función por cada submódulo del RTL de la
Parte C, así después cada testbench compara contra su función:

\begin{center}\small
\begin{tabular}{@{}lll@{}}
\toprule
RTL & modelo & etapa \\
\midrule
\texttt{golay\_syndrome}    & \texttt{syndrome}   & 1 \\
\texttt{golay\_mult\_b}     & \texttt{mult\_b}     & 2 \\
\texttt{popcount12}         & \texttt{popcount12} & 2 \\
\texttt{golay\_row\_search} & \texttt{row\_search} & 2 \\
\texttt{golay\_err\_gen}    & \texttt{err\_gen}    & 3 \\
\texttt{golay\_correct}     & \texttt{correct}    & 3 \\
\bottomrule
\end{tabular}
\end{center}

\texttt{decode()} llama a esas funciones en el orden del pipeline.
Como \texttt{golay\_correct} solo recibe \texttt{i\_rx} e \texttt{i\_err},
\texttt{o\_corrected} es simplemente $e\neq 0$. Cuando
\texttt{row\_search} no encuentra fila, \texttt{o\_idx} y \texttt{o\_res}
no importan.

\subsection{Tabla de síndromes}

La tabla se arma recorriendo todos los errores de peso $\le 4$, calculando
su síndrome con $H$ y guardando, para cada síndrome, los patrones de peso
mínimo:

\begin{center}\small
% generado por gen_parte_b.py -- no editar a mano
\begin{tabular}{@{}cccl@{}}
\toprule
peso del líder & líderes & cosets & \\
\midrule
0 & 1 & 1 & corregible \\
1 & 1 & 24 & corregible \\
2 & 1 & 276 & corregible \\
3 & 1 & 2024 & corregible \\
4 & 6 & 1771 & no corregible \\
\midrule
& & 4096 & \\
\bottomrule
\end{tabular}

\end{center}

Los \nCosetsCorregibles{} síndromes con un solo líder se corrigen. Los
\nCosetsAmbiguos{} restantes tienen seis líderes de peso 4 y no hay forma
de elegir uno, así que se declaran no corregibles. El algoritmo de cuatro
casos no necesita esta tabla, así que se usa como segundo decodificador
para comparar resultados, y los dos dan lo mismo.

\section{Ejercicio 7 --- Caracterización}

\subsection{Errores de peso $\le 4$}

Se decodifican todos los patrones de error de peso 0 a 4 sobre la palabra
nula (alcanza con esa porque el síndrome solo depende de $e$):

\begin{center}\small
% generado por gen_parte_b.py -- no editar a mano
\begin{tabular}{@{}l*{5}{r}@{}}
\toprule
peso de $e$ & 0 & 1 & 2 & 3 & 4 \\
\midrule
patrones & 1 & 24 & 276 & 2024 & 10626 \\
corregidos de forma exacta & 1 & 24 & 276 & 2024 & 0 \\
decodificados a otra palabra & 0 & 0 & 0 & 0 & 0 \\
detectados como no corregibles & 0 & 0 & 0 & 0 & 10626 \\
\bottomrule
\end{tabular}

\end{center}

Da igual que la tabla del enunciado. Nunca se decodifica a una palabra
equivocada: o corrige bien o marca que no puede.

\subsection{Estímulo para la Parte C}

No se generan archivos de vectores: los testbenches de cocotb importan
\texttt{golay/stimulus.py} y generan el estímulo ahí mismo, como en la
Unidad 3. Cada generador devuelve diccionarios con los nombres de los
puertos del módulo:

\begin{center}\small
\begin{tabular}{@{}l@{}}
\toprule
\texttt{for vec in mult\_b\_vectors():} \\
\texttt{\quad dut.i\_vec.value = vec[\string"i\_vec\string"]} \\
\texttt{\quad await Timer(2, \string"ns\string")} \\
\texttt{\quad assert int(dut.o\_vec.value) == vec[\string"o\_vec\string"]} \\
\bottomrule
\end{tabular}
\end{center}

Los módulos de 12 bits se prueban con las 4096 entradas. El decodificador
completo tiene $2^{24}$ entradas posibles, así que se usa una muestra:
las 4096 palabras código, todos los errores de peso 1 a 3 sobre cinco
palabras, y parte de los de peso 4.

\begin{center}\small
% generado por gen_parte_b.py -- no editar a mano
\begin{tabular}{@{}lllr@{}}
\toprule
módulo & generador & casos & criterio \\
\midrule
\texttt{golay\_mult\_b} & \texttt{mult\_b\_vectors} & 4096 & barrido exhaustivo \\
\texttt{popcount12} & \texttt{popcount12\_vectors} & 4096 & barrido exhaustivo \\
\texttt{golay\_row\_search} & \texttt{row\_search\_vectors} & 4096 & barrido exhaustivo \\
\texttt{golay\_encoder} & \texttt{encoder\_vectors} & 4096 & barrido exhaustivo \\
\texttt{golay\_syndrome} & \texttt{syndrome\_vectors} & 6420 & 4096 palabras código $+$ errores \\
\texttt{golay\_decoder} & \texttt{decoder\_vectors} & 15826 & curado, cinco ramas \\
\bottomrule
\end{tabular}

\qquad
% generado por gen_parte_b.py -- no editar a mano
\begin{tabular}{@{}crl@{}}
\toprule
caso & vectores & condición \\
\midrule
1 & 5586 & $w(s)\le 3$ \\
2 & 4740 & $\exists i: w(s\oplus b_i)\le 2$ \\
3 & 1430 & $w(q)\le 3$ \\
4 & 3960 & $\exists i: w(q\oplus b_i)\le 2$ \\
5 & 110 & no corregible \\
\bottomrule
\end{tabular}

\end{center}

\section{Preguntas}

\subsection{¿Por qué el síndrome depende solo del error?}

Con $r=v+e$:
\[
  s = r\,H^T = v\,H^T + e\,H^T = e\,H^T ,
\]
porque $v=m\,G$ y $G\,H^T=0$. El mensaje desaparece de la cuenta.

\subsection{¿Por qué $d_{\min}=8$ detecta 4 errores pero no los corrige?}

\textbf{Detecta:} si $e$ pesa 4, $v+e$ no puede ser otra palabra código,
porque la diferencia entre dos palabras código pesa al menos 8. Entonces
el síndrome no es cero y el error se nota.

\textbf{No corrige:} para corregir $t$ errores hace falta
$d_{\min}\ge 2t+1$, y con $d_{\min}=8$ eso da $t=3$. En la tabla de
síndromes se ve por qué: un error de peso 4 cae en un síndrome que
comparten seis patrones de peso 4, todos igual de probables, y no hay
forma de saber cuál fue. Por eso el decodificador llega al caso 5 y
levanta \texttt{o\_uncorrectable}, que es lo que pasó con $r_3$ en la
Parte A.

\end{document}
